\section{晚明的区域比较与制度实践}

同一项制度在不同地区从不以同一速度、同一成本或同一社会后果展开。账本中的诏令、军报、赋役记录、工程奏报与使节文书提供跨区域比较的起点；地方志、契约、碑刻、遗址与异域材料则提醒我们，中央分类无法覆盖地方实践。本节以事件链而非单一结论呈现制度、文化、环境与边地关系的差异。

\subsection{征解、流民与地方财政}

\paragraph{1635（崇祯八年）三月：叛军占领凤阳} 这一事件使区域差异进入可核查的编年。账本确认的范围是编年候选的年序可由官修与后出材料交叉；具体月日、参与者、战果或数字必须回查所列材料，不能将单一叙事当作定论。，但各材料服务于不同的行政、叙事或保存目的。事件之所以在该时点发生，与；可见的后续变化是它影响后续人事与政策议程；动机和派系标签不能由单一叙事裁定。比较不同地区时，不能把中心政策、地方报数、物质遗存与社会经验混为同一种证据。译写候选以编年材料定位；实录记载反映朝廷选择，崇祯期须另以奏疏、地方及敌对政权材料交叉。

\paragraph{1635（崇祯八年）八月：明军在甘肃被叛军击败} 这一事件使区域差异进入可核查的编年。账本确认的范围是编年候选的年序可由官修与后出材料交叉；具体月日、参与者、战果或数字必须回查所列材料，不能将单一叙事当作定论。，但各材料服务于不同的行政、叙事或保存目的。事件之所以在该时点发生，与；可见的后续变化是它改变了后续调兵、补给或地方秩序的条件；战果和伤亡须分开核对。比较不同地区时，不能把中心政策、地方报数、物质遗存与社会经验混为同一种证据。译写候选以编年材料定位；实录记载反映朝廷选择，崇祯期须另以奏疏、地方及敌对政权材料交叉。

\paragraph{1635（崇祯八年）九月：李自成山西起义军} 这一事件使区域差异进入可核查的编年。账本确认的范围是编年候选的年序可由官修与后出材料交叉；具体月日、参与者、战果或数字必须回查所列材料，不能将单一叙事当作定论。，但各材料服务于不同的行政、叙事或保存目的。事件之所以在该时点发生，与；可见的后续变化是它影响后续人事与政策议程；动机和派系标签不能由单一叙事裁定。比较不同地区时，不能把中心政策、地方报数、物质遗存与社会经验混为同一种证据。译写候选以编年材料定位；实录记载反映朝廷选择，崇祯期须另以奏疏、地方及敌对政权材料交叉。

\paragraph{1635（崇祯八年）：明代使用望远镜来瞄准火炮} 这一事件使区域差异进入可核查的编年。账本确认的范围是编年候选的年序可由官修与后出材料交叉；具体月日、参与者、战果或数字必须回查所列材料，不能将单一叙事当作定论。，但各材料服务于不同的行政、叙事或保存目的。事件之所以在该时点发生，与；可见的后续变化是它提供制度和传播史的锚点，不能据此推断社会整体接受。比较不同地区时，不能把中心政策、地方报数、物质遗存与社会经验混为同一种证据。译写候选以编年材料定位；实录记载反映朝廷选择，崇祯期须另以奏疏、地方及敌对政权材料交叉。

\paragraph{1635（崇祯八年）：辽饷、剿饷、练饷及地方征解的本年文书显示财政负担，定额、征收与实际流向不得混为一谈} 这一事件使区域差异进入可核查的编年。账本确认的范围是来源导向的年度桥接条目：本年材料可定位相关政务线索；具体月日、发文机关、地域和执行结果仍须逐项回查，不能以制度文本或奏报替代实际效果。，但各材料服务于不同的行政、叙事或保存目的。事件之所以在该时点发生，与；可见的后续变化是其法定安排不等于地方执行，后续须以地域材料追踪负担与成效。比较不同地区时，不能把中心政策、地方报数、物质遗存与社会经验混为同一种证据。桥接条目只标示可回查的年度问题与材料链；实录有中央选择性，崇祯期尤其需要奏疏、地方、朝鲜及满文材料交叉。

\paragraph{1635（崇祯八年）：辽东防务、后金（清）军事行动和流动作战集团的本年记录标记多战场互动，军数和胜败须保留分歧} 这一事件使区域差异进入可核查的编年。账本确认的范围是来源导向的年度桥接条目：本年材料可定位相关政务线索；具体月日、发文机关、地域和执行结果仍须逐项回查，不能以制度文本或奏报替代实际效果。，但各材料服务于不同的行政、叙事或保存目的。事件之所以在该时点发生，与；可见的后续变化是它改变了后续调兵、补给或地方秩序的条件；战果和伤亡须分开核对。比较不同地区时，不能把中心政策、地方报数、物质遗存与社会经验混为同一种证据。桥接条目只标示可回查的年度问题与材料链；实录有中央选择性，崇祯期尤其需要奏疏、地方、朝鲜及满文材料交叉。

\subsection{辽东、东南与边地中介}

\paragraph{1635（崇祯八年）：地方结社、宗教、出版或士人文集的本年线索提供战乱中的社会观察，幸存者记忆不能概括全国} 这一事件使区域差异进入可核查的编年。账本确认的范围是来源导向的年度桥接条目：本年材料可定位相关政务线索；具体月日、发文机关、地域和执行结果仍须逐项回查，不能以制度文本或奏报替代实际效果。，但各材料服务于不同的行政、叙事或保存目的。事件之所以在该时点发生，与；可见的后续变化是它提供制度和传播史的锚点，不能据此推断社会整体接受。比较不同地区时，不能把中心政策、地方报数、物质遗存与社会经验混为同一种证据。桥接条目只标示可回查的年度问题与材料链；实录有中央选择性，崇祯期尤其需要奏疏、地方、朝鲜及满文材料交叉。

\paragraph{1636（崇祯九年）：皇太极宣布清朝成立} 这一事件使区域差异进入可核查的编年。账本确认的范围是编年候选的年序可由官修与后出材料交叉；具体月日、参与者、战果或数字必须回查所列材料，不能将单一叙事当作定论。，但各材料服务于不同的行政、叙事或保存目的。事件之所以在该时点发生，与；可见的后续变化是它影响后续人事与政策议程；动机和派系标签不能由单一叙事裁定。比较不同地区时，不能把中心政策、地方报数、物质遗存与社会经验混为同一种证据。译写候选以编年材料定位；实录记载反映朝廷选择，崇祯期须另以奏疏、地方及敌对政权材料交叉。

\paragraph{1636（崇祯九年）：旱灾、蝗灾、疫病、河患与赈济的本年材料连接环境和社会压力，死亡与流离数字须谨慎处理。（材料核查重点：诏令或奏议的形成与发受链）} 这一事件使区域差异进入可核查的编年。账本确认的范围是来源导向的年度桥接条目：本年材料可定位相关政务线索；具体月日、发文机关、地域和执行结果仍须逐项回查，不能以制度文本或奏报替代实际效果。，但各材料服务于不同的行政、叙事或保存目的。事件之所以在该时点发生，与；可见的后续变化是赈济、迁徙和损失的实际范围仍须以受灾地区材料分别核验。比较不同地区时，不能把中心政策、地方报数、物质遗存与社会经验混为同一种证据。桥接条目只标示可回查的年度问题与材料链；实录有中央选择性，崇祯期尤其需要奏疏、地方、朝鲜及满文材料交叉。；本条特别核对诏令或奏议的形成与发受链。

\paragraph{1636（崇祯九年）：辽东防务、后金（清）军事行动和流动作战集团的本年记录标记多战场互动，军数和胜败须保留分歧。（材料核查重点：诏令或奏议的形成与发受链）} 这一事件使区域差异进入可核查的编年。账本确认的范围是来源导向的年度桥接条目：本年材料可定位相关政务线索；具体月日、发文机关、地域和执行结果仍须逐项回查，不能以制度文本或奏报替代实际效果。，但各材料服务于不同的行政、叙事或保存目的。事件之所以在该时点发生，与；可见的后续变化是它改变了后续调兵、补给或地方秩序的条件；战果和伤亡须分开核对。比较不同地区时，不能把中心政策、地方报数、物质遗存与社会经验混为同一种证据。桥接条目只标示可回查的年度问题与材料链；实录有中央选择性，崇祯期尤其需要奏疏、地方、朝鲜及满文材料交叉。；本条特别核对诏令或奏议的形成与发受链。

\paragraph{1636（崇祯九年）：朝鲜、辽东和海上关系的本年材料提供外部观察位置，明、后金（清）与地方势力的称谓需具体化。（材料核查重点：诏令或奏议的形成与发受链）} 这一事件使区域差异进入可核查的编年。账本确认的范围是来源导向的年度桥接条目：本年材料可定位相关政务线索；具体月日、发文机关、地域和执行结果仍须逐项回查，不能以制度文本或奏报替代实际效果。，但各材料服务于不同的行政、叙事或保存目的。事件之所以在该时点发生，与；可见的后续变化是它为后续册封、互市或海防安排留下文书与跨政权比较线索。比较不同地区时，不能把中心政策、地方报数、物质遗存与社会经验混为同一种证据。桥接条目只标示可回查的年度问题与材料链；实录有中央选择性，崇祯期尤其需要奏疏、地方、朝鲜及满文材料交叉。；本条特别核对诏令或奏议的形成与发受链。

\paragraph{1636（崇祯九年）：辽东防务、后金（清）军事行动和流动作战集团的本年记录标记多战场互动，军数和胜败须保留分歧。（材料核查重点：报称信息的来源、抄传与行政分类）} 这一事件使区域差异进入可核查的编年。账本确认的范围是来源导向的年度桥接条目：本年材料可定位相关政务线索；具体月日、发文机关、地域和执行结果仍须逐项回查，不能以制度文本或奏报替代实际效果。，但各材料服务于不同的行政、叙事或保存目的。事件之所以在该时点发生，与；可见的后续变化是它改变了后续调兵、补给或地方秩序的条件；战果和伤亡须分开核对。比较不同地区时，不能把中心政策、地方报数、物质遗存与社会经验混为同一种证据。桥接条目只标示可回查的年度问题与材料链；实录有中央选择性，崇祯期尤其需要奏疏、地方、朝鲜及满文材料交叉。；本条特别核对报称信息的来源、抄传与行政分类。

\subsection{城市、道路与文化资源}

\paragraph{1636（崇祯九年）：朝鲜、辽东和海上关系的本年材料提供外部观察位置，明、后金（清）与地方势力的称谓需具体化。（材料核查重点：报称信息的来源、抄传与行政分类）} 这一事件使区域差异进入可核查的编年。账本确认的范围是来源导向的年度桥接条目：本年材料可定位相关政务线索；具体月日、发文机关、地域和执行结果仍须逐项回查，不能以制度文本或奏报替代实际效果。，但各材料服务于不同的行政、叙事或保存目的。事件之所以在该时点发生，与；可见的后续变化是它为后续册封、互市或海防安排留下文书与跨政权比较线索。比较不同地区时，不能把中心政策、地方报数、物质遗存与社会经验混为同一种证据。桥接条目只标示可回查的年度问题与材料链；实录有中央选择性，崇祯期尤其需要奏疏、地方、朝鲜及满文材料交叉。；本条特别核对报称信息的来源、抄传与行政分类。

\paragraph{1636（崇祯九年）：旱灾、蝗灾、疫病、河患与赈济的本年材料连接环境和社会压力，死亡与流离数字须谨慎处理。（材料核查重点：报称信息的来源、抄传与行政分类）} 这一事件使区域差异进入可核查的编年。账本确认的范围是来源导向的年度桥接条目：本年材料可定位相关政务线索；具体月日、发文机关、地域和执行结果仍须逐项回查，不能以制度文本或奏报替代实际效果。，但各材料服务于不同的行政、叙事或保存目的。事件之所以在该时点发生，与；可见的后续变化是赈济、迁徙和损失的实际范围仍须以受灾地区材料分别核验。比较不同地区时，不能把中心政策、地方报数、物质遗存与社会经验混为同一种证据。桥接条目只标示可回查的年度问题与材料链；实录有中央选择性，崇祯期尤其需要奏疏、地方、朝鲜及满文材料交叉。；本条特别核对报称信息的来源、抄传与行政分类。

\paragraph{1637（崇祯十年）：天启末至崇祯朝的任免、奏疏和廷议构成本年中枢危机线索；崇祯无实录，必须以多类材料互校。（材料核查重点：诏令或奏议的形成与发受链）} 这一事件使区域差异进入可核查的编年。账本确认的范围是来源导向的年度桥接条目：本年材料可定位相关政务线索；具体月日、发文机关、地域和执行结果仍须逐项回查，不能以制度文本或奏报替代实际效果。，但各材料服务于不同的行政、叙事或保存目的。事件之所以在该时点发生，与；可见的后续变化是它影响后续人事与政策议程；动机和派系标签不能由单一叙事裁定。比较不同地区时，不能把中心政策、地方报数、物质遗存与社会经验混为同一种证据。桥接条目只标示可回查的年度问题与材料链；实录有中央选择性，崇祯期尤其需要奏疏、地方、朝鲜及满文材料交叉。；本条特别核对诏令或奏议的形成与发受链。

\paragraph{1637（崇祯十年）：天启末至崇祯朝的任免、奏疏和廷议构成本年中枢危机线索；崇祯无实录，必须以多类材料互校。（材料核查重点：报称信息的来源、抄传与行政分类）} 这一事件使区域差异进入可核查的编年。账本确认的范围是来源导向的年度桥接条目：本年材料可定位相关政务线索；具体月日、发文机关、地域和执行结果仍须逐项回查，不能以制度文本或奏报替代实际效果。，但各材料服务于不同的行政、叙事或保存目的。事件之所以在该时点发生，与；可见的后续变化是它影响后续人事与政策议程；动机和派系标签不能由单一叙事裁定。比较不同地区时，不能把中心政策、地方报数、物质遗存与社会经验混为同一种证据。桥接条目只标示可回查的年度问题与材料链；实录有中央选择性，崇祯期尤其需要奏疏、地方、朝鲜及满文材料交叉。；本条特别核对报称信息的来源、抄传与行政分类。

\paragraph{1637（崇祯十年）：朝鲜、辽东和海上关系的本年材料提供外部观察位置，明、后金（清）与地方势力的称谓需具体化。（材料核查重点：诏令或奏议的形成与发受链）} 这一事件使区域差异进入可核查的编年。账本确认的范围是来源导向的年度桥接条目：本年材料可定位相关政务线索；具体月日、发文机关、地域和执行结果仍须逐项回查，不能以制度文本或奏报替代实际效果。，但各材料服务于不同的行政、叙事或保存目的。事件之所以在该时点发生，与；可见的后续变化是它为后续册封、互市或海防安排留下文书与跨政权比较线索。比较不同地区时，不能把中心政策、地方报数、物质遗存与社会经验混为同一种证据。桥接条目只标示可回查的年度问题与材料链；实录有中央选择性，崇祯期尤其需要奏疏、地方、朝鲜及满文材料交叉。；本条特别核对诏令或奏议的形成与发受链。

\paragraph{1637（崇祯十年）：地方结社、宗教、出版或士人文集的本年线索提供战乱中的社会观察，幸存者记忆不能概括全国。（材料核查重点：诏令或奏议的形成与发受链）} 这一事件使区域差异进入可核查的编年。账本确认的范围是来源导向的年度桥接条目：本年材料可定位相关政务线索；具体月日、发文机关、地域和执行结果仍须逐项回查，不能以制度文本或奏报替代实际效果。，但各材料服务于不同的行政、叙事或保存目的。事件之所以在该时点发生，与；可见的后续变化是它提供制度和传播史的锚点，不能据此推断社会整体接受。比较不同地区时，不能把中心政策、地方报数、物质遗存与社会经验混为同一种证据。桥接条目只标示可回查的年度问题与材料链；实录有中央选择性，崇祯期尤其需要奏疏、地方、朝鲜及满文材料交叉。；本条特别核对诏令或奏议的形成与发受链。

\subsection{环境、工程与地方应对}

\paragraph{1637（崇祯十年）：朝鲜、辽东和海上关系的本年材料提供外部观察位置，明、后金（清）与地方势力的称谓需具体化。（材料核查重点：报称信息的来源、抄传与行政分类）} 这一事件使区域差异进入可核查的编年。账本确认的范围是来源导向的年度桥接条目：本年材料可定位相关政务线索；具体月日、发文机关、地域和执行结果仍须逐项回查，不能以制度文本或奏报替代实际效果。，但各材料服务于不同的行政、叙事或保存目的。事件之所以在该时点发生，与；可见的后续变化是它为后续册封、互市或海防安排留下文书与跨政权比较线索。比较不同地区时，不能把中心政策、地方报数、物质遗存与社会经验混为同一种证据。桥接条目只标示可回查的年度问题与材料链；实录有中央选择性，崇祯期尤其需要奏疏、地方、朝鲜及满文材料交叉。；本条特别核对报称信息的来源、抄传与行政分类。

\paragraph{1637（崇祯十年）：地方结社、宗教、出版或士人文集的本年线索提供战乱中的社会观察，幸存者记忆不能概括全国。（材料核查重点：报称信息的来源、抄传与行政分类）} 这一事件使区域差异进入可核查的编年。账本确认的范围是来源导向的年度桥接条目：本年材料可定位相关政务线索；具体月日、发文机关、地域和执行结果仍须逐项回查，不能以制度文本或奏报替代实际效果。，但各材料服务于不同的行政、叙事或保存目的。事件之所以在该时点发生，与；可见的后续变化是它提供制度和传播史的锚点，不能据此推断社会整体接受。比较不同地区时，不能把中心政策、地方报数、物质遗存与社会经验混为同一种证据。桥接条目只标示可回查的年度问题与材料链；实录有中央选择性，崇祯期尤其需要奏疏、地方、朝鲜及满文材料交叉。；本条特别核对报称信息的来源、抄传与行政分类。

\paragraph{1637（崇祯十年）：天启末至崇祯朝的任免、奏疏和廷议构成本年中枢危机线索；崇祯无实录，必须以多类材料互校。（材料核查重点：同年地方材料所见的执行范围）} 这一事件使区域差异进入可核查的编年。账本确认的范围是来源导向的年度桥接条目：本年材料可定位相关政务线索；具体月日、发文机关、地域和执行结果仍须逐项回查，不能以制度文本或奏报替代实际效果。，但各材料服务于不同的行政、叙事或保存目的。事件之所以在该时点发生，与；可见的后续变化是它影响后续人事与政策议程；动机和派系标签不能由单一叙事裁定。比较不同地区时，不能把中心政策、地方报数、物质遗存与社会经验混为同一种证据。桥接条目只标示可回查的年度问题与材料链；实录有中央选择性，崇祯期尤其需要奏疏、地方、朝鲜及满文材料交叉。；本条特别核对同年地方材料所见的执行范围。

\paragraph{1638（崇祯十一年）：清朝征服山东} 这一事件使区域差异进入可核查的编年。账本确认的范围是编年候选的年序可由官修与后出材料交叉；具体月日、参与者、战果或数字必须回查所列材料，不能将单一叙事当作定论。，但各材料服务于不同的行政、叙事或保存目的。事件之所以在该时点发生，与；可见的后续变化是它改变了后续调兵、补给或地方秩序的条件；战果和伤亡须分开核对。比较不同地区时，不能把中心政策、地方报数、物质遗存与社会经验混为同一种证据。译写候选以编年材料定位；实录记载反映朝廷选择，崇祯期须另以奏疏、地方及敌对政权材料交叉。

\paragraph{1638（崇祯十一年）：明军在晋豫边境大败} 这一事件使区域差异进入可核查的编年。账本确认的范围是编年候选的年序可由官修与后出材料交叉；具体月日、参与者、战果或数字必须回查所列材料，不能将单一叙事当作定论。，但各材料服务于不同的行政、叙事或保存目的。事件之所以在该时点发生，与；可见的后续变化是它改变了后续调兵、补给或地方秩序的条件；战果和伤亡须分开核对。比较不同地区时，不能把中心政策、地方报数、物质遗存与社会经验混为同一种证据。译写候选以编年材料定位；实录记载反映朝廷选择，崇祯期须另以奏疏、地方及敌对政权材料交叉。

\paragraph{1638（崇祯十一年）：地方结社、宗教、出版或士人文集的本年线索提供战乱中的社会观察，幸存者记忆不能概括全国。（材料核查重点：诏令或奏议的形成与发受链）} 这一事件使区域差异进入可核查的编年。账本确认的范围是来源导向的年度桥接条目：本年材料可定位相关政务线索；具体月日、发文机关、地域和执行结果仍须逐项回查，不能以制度文本或奏报替代实际效果。，但各材料服务于不同的行政、叙事或保存目的。事件之所以在该时点发生，与；可见的后续变化是它提供制度和传播史的锚点，不能据此推断社会整体接受。比较不同地区时，不能把中心政策、地方报数、物质遗存与社会经验混为同一种证据。桥接条目只标示可回查的年度问题与材料链；实录有中央选择性，崇祯期尤其需要奏疏、地方、朝鲜及满文材料交叉。；本条特别核对诏令或奏议的形成与发受链。

\subsection{环境、工程与地方应对}

\paragraph{1638（崇祯十一年）：旱灾、蝗灾、疫病、河患与赈济的本年材料连接环境和社会压力，死亡与流离数字须谨慎处理。（材料核查重点：诏令或奏议的形成与发受链）} 这一事件使区域差异进入可核查的编年。账本确认的范围是来源导向的年度桥接条目：本年材料可定位相关政务线索；具体月日、发文机关、地域和执行结果仍须逐项回查，不能以制度文本或奏报替代实际效果。，但各材料服务于不同的行政、叙事或保存目的。事件之所以在该时点发生，与；可见的后续变化是赈济、迁徙和损失的实际范围仍须以受灾地区材料分别核验。比较不同地区时，不能把中心政策、地方报数、物质遗存与社会经验混为同一种证据。桥接条目只标示可回查的年度问题与材料链；实录有中央选择性，崇祯期尤其需要奏疏、地方、朝鲜及满文材料交叉。；本条特别核对诏令或奏议的形成与发受链。

\paragraph{1638（崇祯十一年）：地方结社、宗教、出版或士人文集的本年线索提供战乱中的社会观察，幸存者记忆不能概括全国。（材料核查重点：报称信息的来源、抄传与行政分类）} 这一事件使区域差异进入可核查的编年。账本确认的范围是来源导向的年度桥接条目：本年材料可定位相关政务线索；具体月日、发文机关、地域和执行结果仍须逐项回查，不能以制度文本或奏报替代实际效果。，但各材料服务于不同的行政、叙事或保存目的。事件之所以在该时点发生，与；可见的后续变化是它提供制度和传播史的锚点，不能据此推断社会整体接受。比较不同地区时，不能把中心政策、地方报数、物质遗存与社会经验混为同一种证据。桥接条目只标示可回查的年度问题与材料链；实录有中央选择性，崇祯期尤其需要奏疏、地方、朝鲜及满文材料交叉。；本条特别核对报称信息的来源、抄传与行政分类。

\paragraph{1638（崇祯十一年）：旱灾、蝗灾、疫病、河患与赈济的本年材料连接环境和社会压力，死亡与流离数字须谨慎处理。（材料核查重点：报称信息的来源、抄传与行政分类）} 这一事件使区域差异进入可核查的编年。账本确认的范围是来源导向的年度桥接条目：本年材料可定位相关政务线索；具体月日、发文机关、地域和执行结果仍须逐项回查，不能以制度文本或奏报替代实际效果。，但各材料服务于不同的行政、叙事或保存目的。事件之所以在该时点发生，与；可见的后续变化是赈济、迁徙和损失的实际范围仍须以受灾地区材料分别核验。比较不同地区时，不能把中心政策、地方报数、物质遗存与社会经验混为同一种证据。桥接条目只标示可回查的年度问题与材料链；实录有中央选择性，崇祯期尤其需要奏疏、地方、朝鲜及满文材料交叉。；本条特别核对报称信息的来源、抄传与行政分类。

\paragraph{1638（崇祯十一年）：辽饷、剿饷、练饷及地方征解的本年文书显示财政负担，定额、征收与实际流向不得混为一谈} 这一事件使区域差异进入可核查的编年。账本确认的范围是来源导向的年度桥接条目：本年材料可定位相关政务线索；具体月日、发文机关、地域和执行结果仍须逐项回查，不能以制度文本或奏报替代实际效果。，但各材料服务于不同的行政、叙事或保存目的。事件之所以在该时点发生，与；可见的后续变化是其法定安排不等于地方执行，后续须以地域材料追踪负担与成效。比较不同地区时，不能把中心政策、地方报数、物质遗存与社会经验混为同一种证据。桥接条目只标示可回查的年度问题与材料链；实录有中央选择性，崇祯期尤其需要奏疏、地方、朝鲜及满文材料交叉。

\paragraph{1639（崇祯十二年）：西班牙人和菲律宾人在吕宋岛屠杀了两万名中国人} 这一事件使区域差异进入可核查的编年。账本确认的范围是编年候选的年序可由官修与后出材料交叉；具体月日、参与者、战果或数字必须回查所列材料，不能将单一叙事当作定论。，但各材料服务于不同的行政、叙事或保存目的。事件之所以在该时点发生，与；可见的后续变化是它为后续册封、互市或海防安排留下文书与跨政权比较线索。比较不同地区时，不能把中心政策、地方报数、物质遗存与社会经验混为同一种证据。译写候选以编年材料定位；实录记载反映朝廷选择，崇祯期须另以奏疏、地方及敌对政权材料交叉。

\paragraph{1639（崇祯十二年）：禁止来自澳门的葡萄牙商人进入长崎} 这一事件使区域差异进入可核查的编年。账本确认的范围是编年候选的年序可由官修与后出材料交叉；具体月日、参与者、战果或数字必须回查所列材料，不能将单一叙事当作定论。，但各材料服务于不同的行政、叙事或保存目的。事件之所以在该时点发生，与；可见的后续变化是它为后续册封、互市或海防安排留下文书与跨政权比较线索。比较不同地区时，不能把中心政策、地方报数、物质遗存与社会经验混为同一种证据。译写候选以编年材料定位；实录记载反映朝廷选择，崇祯期须另以奏疏、地方及敌对政权材料交叉。
